% ============================================================
% Chapter 2 -- Scientific Background and Related Work
% ============================================================

\section{Retrieval-Augmented Generation}

A RAG system has three stages \citep{lewis2020rag,gao2023ragsurvey}. An \emph{indexer} embeds the corpus offline. A \emph{retriever} returns the documents most similar to the query. A \emph{generator} LLM produces the answer conditioned on them. ``Modular'' RAG architectures add pre-retrieval steps (query rewriting) and post-retrieval steps (re-ranking, filtering, compression) around the retriever \citep{gao2023ragsurvey}. These additions exist because the generator is sensitive to what it receives. Irrelevant passages distract it \citep{shi2023distractors}, and relevant information placed in the middle of a long context is often ignored \citep{liu2024lost}. Black-box generators can be improved by tuning the retriever alone \citep{shi2024replug}, which confirms that retrieval is a lever on end-to-end quality. In medicine, the MIRAGE benchmark and the MedRAG toolkit showed that the choice of corpus, retriever and number of retrieved snippets changes answer accuracy substantially \citep{xiong2024medrag}. BERGEN provides a comparable benchmarking library for general-domain RAG \citep{rau2024bergen}. RAG outputs are evaluated by humans, by model-based judges \citep{zheng2023judging,es2024ragas}, or component by component. We follow the component-level view and measure retrieval directly, because retrieval errors bound everything downstream.

\section{Dense Retrieval and Nearest-Neighbour Search}

Dense retrievers encode queries and documents independently with a bi-encoder and score them by inner product or cosine similarity \citep{karpukhin2020dpr,reimers2019sbert}. Document vectors are pre-computed and stored in a vector database. At query time, a graph-based index such as HNSW \citep{malkov2020hnsw} finds approximate nearest neighbours by greedy best-first traversal from entry points in a proximity graph. Strong general-purpose encoders include MPNet-based sentence transformers \citep{song2020mpnet}, which we use, and Contriever \citep{izacard2022contriever}. Biomedical encoders are pre-trained on domain text \citep{lee2020biobert,gu2021pubmedbert} or trained contrastively on PubMed search logs \citep{jin2023medcpt}. Dense models degrade outside their training distribution, and BEIR shows large out-of-domain gaps, including on biomedical tasks \citep{thakur2021beir}. Sparse lexical retrieval with BM25 \citep{robertson2009bm25} remains a strong and complementary baseline.

\section{Geometry of Embedding Spaces and the Query--Document Gap}

Contextual embeddings are \emph{anisotropic}: rather than spreading over the hypersphere, they occupy a narrow cone, so random pairs have high cosine similarity \citep{ethayarajh2019contextual}. Contrastively trained two-tower models can additionally exhibit a \emph{modality gap}: the two input types are embedded in separate regions of the shared space, and the gap is created by initialisation and preserved by the contrastive objective \citep{liang2022modalitygap}. Questions and documents differ in length, style and vocabulary, so a bi-encoder applied to question answering faces an analogous \emph{query--document gap}. In that case the query may lie far from the dense region of documents that answer it. Work on hard negatives shows that documents close to the query in embedding space are often not relevant \citep{xiong2021ance}, and topic-aware training improves robustness to topically diverse queries \citep{hofstatter2021tasb}. Our preliminary results (Chapter~\ref{chapter:preliminary}) measure this gap directly in a biomedical corpus.

\section{Re-ranking and LLM Relevance Judgements}

Cross-encoders read the query and the document jointly and re-rank a candidate list far more accurately than bi-encoders, at a much higher cost per document \citep{nogueira2019passage}. LLMs can be used as zero-shot re-rankers, pointwise or listwise \citep{sun2023rankgpt}, and as relevance assessors whose labels agree closely with human judgements \citep{thomas2024llmjudge,faggioli2023perspectives}. The broader LLM-as-a-judge literature documents both the reliability and the biases of LLM evaluators \citep{zheng2023judging}. Two limitations matter here. First, cost: every judged document is an LLM call. Second, resolution: pointwise judges asked for a label on a small integer scale assign the same label to many documents. \citet{zhuang2024beyond} show that finer-grained relevance labels and expected-relevance scores computed from token probabilities improve zero-shot LLM rankers. This is directly relevant to the tie problem we observe.

\section{Query Reformulation and Relevance Feedback}

Relevance feedback moves the query representation towards documents marked relevant and away from documents marked non-relevant. The Rocchio formula \citep{rocchio1971relevance} has exactly this positive and negative form. Pseudo-relevance feedback (PRF) assumes that the top-ranked documents are relevant. ANCE-PRF learns a new query embedding from the query and its top-$k$ dense-retrieval results \citep{yu2021anceprf}. LLM-based reformulation generates text in the document's register instead: HyDE embeds an LLM-generated hypothetical answer document and retrieves its neighbours \citep{gao2023hyde}, and Query2doc appends a generated pseudo-document to the query \citep{wang2023query2doc}. These methods change the \emph{anchor} of retrieval. Our setting differs in that true relevance feedback, from the judge, arrives one document at a time and has a cost. The question is therefore how to choose which document to obtain feedback on.

\section{Adaptive RAG and Retrieval-Failure Prediction}

Several systems decide adaptively whether and how to retrieve. FLARE retrieves when the generator's confidence drops \citep{jiang2023flare}. Self-RAG trains the generator to emit reflection tokens that critique retrieval \citep{asai2024selfrag}. Corrective RAG (CRAG) uses a light retrieval evaluator to detect poor retrieval and trigger corrective actions such as web search \citep{yan2024crag}. Query performance prediction (QPP) predicts retrieval quality before or after retrieval from properties of the query and of the score distribution \citep{cronentownsend2002predicting,carmel2010estimating}. These ideas motivate our gated policies: the expensive search should run only on queries where the cheap retriever is likely to have failed.

\section{Search with an Expensive Evaluation Function}

Our problem is an instance of a classical pattern: optimising an expensive black-box function with the help of a cheap surrogate under a budget. In heuristic search, best-first search expands nodes in the order of an evaluation function $f$ \citep{pearl1984heuristics}. A dense retriever is best-first search with $f=\dist(q,\cdot)$, and the number of expansions until the goal is reached is the rank of the target. Bayesian optimisation chooses the next point to evaluate by trading off exploitation and exploration of a surrogate model \citep{shahriari2016bayesopt}. Active learning \citep{settles2009active} and preference elicitation in recommender systems \citep{rubens2015active} choose the next query to an expensive oracle (a user or an annotator) to maximise information. Semantle-style games and trilateration localise a hidden point from distance feedback to probe points. Our contribution is to bring this view to RAG, with the LLM judge as the oracle and the embedding geometry as the surrogate.

\section{Evaluation and Statistical Methodology}

Retrieval experiments are usually compared on the same set of queries, so paired significance tests and an explicit analysis of statistical power are needed to avoid both false discoveries and under-powered conclusions \citep{smucker2007comparison,sakai2016statistical}. This is particularly important here, because the queries on which the retriever fails, where improvements can occur, are a minority of each dataset. Candidate benchmarks include PubMedQA \citep{jin2019pubmedqa}, BioASQ \citep{tsatsaronis2015bioasq,nentidis2023bioasq11}, MS~MARCO \citep{nguyen2016msmarco} and Natural Questions \citep{kwiatkowski2019nq}. Multiple-choice benchmarks such as MedQA \citep{jin2021medqa} and MMLU \citep{hendrycks2021mmlu} are relevant for end-to-end evaluation but have no document-level ground truth.

\section{Summary and Research Gap}

Existing work either improves the cheap signal (better encoders, query rewriting, PRF), applies the expensive signal to a fixed candidate list (cross-encoder and LLM re-ranking), or decides \emph{whether} to retrieve (adaptive RAG). To our knowledge, little work treats the LLM judge as an \emph{expensive oracle inside a sequential search over the embedding space}, in which each judgement updates where to look next and the number of judge calls is the explicit cost. Little work has also related the success of such a search to measurable geometric properties of the query--document space. This proposal addresses both.
